\section{Nonlinear Model, Trim, and Linearization}
\label{sec:model}

\subsection{Nonlinear Model}
The aircraft is a small hybrid VTOL platform with a conventional
fixed-wing configuration for cruise. Its physical parameters are
$m=2.2$\,kg, wing area $S=0.27$\,m$^2$, span $b=1.2$\,m, mean aerodynamic
chord $\bar{c}=0.226$\,m (aspect ratio $AR=5.33$, Oswald factor $e=0.8$),
and inertia tensor (kg\,m$^2$)
\begin{equation}
\mathbf{I}=\begin{bmatrix}
0.1441 & 0 & -0.00167\\
0 & 0.1155 & 0\\
-0.00167 & 0 & 0.2572
\end{bmatrix}.
\end{equation}

The complete six-degree-of-freedom nonlinear model of this
aircraft~\cite{ferreira2019modeling,repo} follows the standard rigid-body
formulation~\cite{stevens2016aircraft} and is implemented in Simulink with a
state vector of 14 states (body velocities, angular rates, Euler angles,
and navigation states). The aerodynamic model uses the usual
coefficient build-up,
\begin{equation}
\begin{aligned}
C_L &= C_{L_0} + C_{L_\alpha}\alpha + C_{L_q}\frac{q\bar{c}}{2V_0}, \quad
C_D = C_{D_0} + \frac{C_L^2}{\pi e AR},\\
C_m &= C_{m_0} + C_{m_\alpha}\alpha + C_{m_{\delta_e}}\delta_e +
C_{m_q}\frac{q\bar{c}}{2V_0},
\end{aligned}
\end{equation}
with analogous expressions for the lateral--directional coefficients
$C_Y$, $C_l$, and $C_n$ in $\beta$, $p$, $r$, $\delta_a$, and $\delta_r$.
The stability and control derivatives were obtained by system
identification from flight-test data of the actual
aircraft~\cite{angelo2026} (longitudinal Model~4 and
lateral--directional Model~2 of that identification campaign),
following the time-domain output-error practice that is standard for
flight-vehicle identification~\cite{jategaonkar2006flight} and for
small UAVs in particular~\cite{dorobantu2012system}. Thrust is
modeled as $T = T_{max}\,\delta_t\,(\rho/\rho_0)^{0.8}$ with
$T_{max}=14$\,N.

In the fixed-wing flight mode considered here the hover rotors are
inactive, so the plant input vector comprises the four aerodynamic
commands and the three components of the wind velocity, the latter
expressed in the local NED frame:
\begin{equation}
\label{eq:inputvec}
U=[\,\underbrace{\delta_t\;\;\delta_e\;\;\delta_a\;\;\delta_r}_{\text{commands}}
\;\;|\;\;
\underbrace{w_N\;\;w_E\;\;w_D}_{\text{wind (NED)}}\,]^{\!\top}.
\end{equation}
The wind components are rotated into body axes and subtracted from the
inertial velocities before the aerodynamic triad is formed, so that
airspeed, angle of attack, and sideslip are computed from the
\emph{air-relative} velocity,
$V_T=\lVert[\,u\;v\;w\,]^{\!\top}-L_{bn}[\,w_N\;w_E\;w_D\,]^{\!\top}\rVert$,
in the usual way~\cite{stevens2016aircraft}. Wind therefore enters the model as a
physical disturbance input to the aerodynamics rather than as an additive
signal on the outputs; it is identically zero at trim and in the nominal
simulations, and is exercised as a disturbance in the robustness study of
Section~\ref{sec:robustness}.

Actuator dynamics are included explicitly, since the controllers are
ultimately deployed on hardware: each control surface is limited to
$\pm25^\circ$ in position and $150^\circ$/s in rate, followed by a
first-order lag of $\tau_{act}=0.05$\,s (20\,rad/s bandwidth); the
throttle channel is limited to $[0,1]$ with a rate limit of 1\,s$^{-1}$
and a spool-up lag of $\tau_{eng}=0.30$\,s. Without these, a reference
step produces an unrealizable impulsive surface deflection
(${\sim}400^\circ$/s in a single step) that no servo --- and no HIL
target --- could reproduce.

\subsection{Trim Condition}
\label{sec:trim}
The aircraft is trimmed numerically --- a Newton solution of
$\dot{x}=f(x,\delta)=0$ under the flight-condition constraints, the
standard formulation of the general aircraft trim
problem~\cite{marco2007trim,stevens2016aircraft} --- in
steady, wings-level horizontal flight at airspeed $\VT^e = 12.0$\,m/s,
altitude $h^e = 600$\,m, and flight-path angle $\gamma^e = 0$. The
resulting equilibrium is stated explicitly below, as it anchors every
simulation and every embedded test in this article.

\emph{State vector at trim} (body axes; unlisted states are zero):
\begin{equation}
\label{eq:xtrim}
\begin{aligned}
u^e &= 11.621~\text{m/s}, \quad
w^e = 2.992~\text{m/s}, \\
\theta^e &= 0.25196~\text{rad} = 14.44^\circ, \quad
h^e = 600~\text{m},
\end{aligned}
\end{equation}
so that $\VT^e=\sqrt{(u^e)^2+(w^e)^2}=12.0$\,m/s and, since
$\alpha^e=\arctan(w^e/u^e)=14.44^\circ=\theta^e$, the flight-path angle is
indeed $\gamma^e=\theta^e-\alpha^e=0$.

\emph{Output vector at trim}:
\begin{equation}
\label{eq:ytrim}
\VT^e = 12.0~\text{m/s},\;\;
\alpha^e = 14.44^\circ,\;\;
\beta^e = 0,\;\;
\gamma^e = 0,
\end{equation}
with $p^e=q^e=r^e=0$, $\phi^e=\psi^e=0$, $\theta^e=14.44^\circ$, and
$h^e=600$\,m.

\emph{Command vector at trim} (control surfaces and throttle):
\begin{equation}
\label{eq:dtrim}
\delta^e =
\begin{bmatrix}
\delta_{t}^{e} \\ \delta_{e}^{e} \\ \delta_{a}^{e} \\ \delta_{r}^{e}
\end{bmatrix}
=
\begin{bmatrix}
0.2836 \\ 0.1330~\text{rad}~(7.62^\circ) \\ 0 \\ 0
\end{bmatrix},
\end{equation}
with the hover-rotor commands identically zero in this flight mode.
Throughout the article, the controllers act on \emph{increments} around
this trim; the trim commands of~\eqref{eq:dtrim} are added at the
controller output (see the remark in Section~\ref{sec:sil-trim}).

\subsection{Linearization}
The nonlinear model is linearized about the trim
condition~\eqref{eq:xtrim}--\eqref{eq:dtrim}, decoupling into the standard
longitudinal and lateral--directional subsystems used for controller
design in Section~\ref{sec:control}. With longitudinal states
$x_{lon}=[\,u\;\alpha\;q\;\theta\;h\,]^{\!\top}$ and inputs
$[\,\delta_t\;\delta_e\,]^{\!\top}$:
\begin{equation}
\resizebox{.91\linewidth}{!}{$
A_{lon}=\begin{bmatrix}
0.124 & 13.29 & -2.643 & -9.497 & 0\\
-0.135 & -2.989 & 0.751 & 0 & 0\\
0 & -38.75 & -4.108 & 0 & 0\\
0 & 0 & 1 & 0 & 0\\
0 & -12 & 0 & 12 & 0
\end{bmatrix}\!,\;
B_{lon}=\begin{bmatrix}
6.071 & 0\\ -0.126 & 0\\ 0 & 45.98\\ 0 & 0\\ 0 & 0
\end{bmatrix}
$}
\end{equation}
and with lateral--directional states
$x_{lat}=[\,\beta\;p\;r\;\phi\;\psi\,]^{\!\top}$ and inputs
$[\,\delta_a\;\delta_r\,]^{\!\top}$:
\begin{equation}
\resizebox{.91\linewidth}{!}{$
A_{lat}=\begin{bmatrix}
-0.478 & 0.220 & -0.967 & 0.791 & 0\\
-51.67 & -12.43 & 9.362 & 0 & 0\\
15.74 & -0.648 & -2.355 & 0 & 0\\
0 & 1 & 0.257 & 0 & 0\\
0 & 0 & 1.033 & 0 & 0
\end{bmatrix}\!,\;
B_{lat}=\begin{bmatrix}
0 & 0.040\\ -132.9 & -12.69\\ 3.282 & 19.09\\ 0 & 0\\ 0 & 0
\end{bmatrix}
$}
\end{equation}
(entries below $10^{-3}$ shown as zero).

\subsection{Open-Loop Stability}
\label{sec:stability}
Table~\ref{tab:poles} lists the eigenvalues of $A_{lon}$ and $A_{lat}$ and
the corresponding classical modes~\cite{etkin1996dynamics,cook2007flight}. The longitudinal dynamics are stable,
with a well-damped short period and a lightly damped phugoid
($\zeta=0.085$). The lateral--directional dynamics, however, contain an
\emph{unstable spiral mode} at $s=+0.233$\,s$^{-1}$ --- an amplitude
doubling time of only $t_2=\ln 2/0.233 \approx 3.0$\,s --- so the aircraft
is open-loop unstable and requires closed-loop stabilization for any
practical operation.

\begin{table}[!t]
\caption{Open-loop modes at the trim condition ($\VT^e=12$ m/s, $h^e=600$ m).}
\label{tab:poles}
\centering
\begin{tabular}{llcc}
\toprule
Mode & Eigenvalues [s$^{-1}$] & $\omega_n$ [rad/s] & $\zeta$ \\
\midrule
Short period & $-3.39 \pm 5.31i$ & 6.30 & 0.54 \\
Phugoid & $-0.094 \pm 1.11i$ & 1.12 & 0.085 \\
Altitude & $-1.7\times10^{-4}$ & --- & --- \\
Roll subsidence & $-11.33$ & --- & --- \\
Dutch roll & $-2.08 \pm 4.59i$ & 5.04 & 0.41 \\
\textbf{Spiral (unstable)} & $\mathbf{+0.233}$ & --- & --- \\
Heading & $0$ & --- & --- \\
\bottomrule
\end{tabular}
\end{table}

This analysis is corroborated by an open-loop simulation of the
\emph{nonlinear} model with the commands frozen at the trim
values~\eqref{eq:dtrim} (Fig.~\ref{fig:openloop}). Initialized exactly at
the trim state~\eqref{eq:xtrim}, the aircraft remains at equilibrium for
the full 300\,s (deviations below $10^{-5}$), confirming the trim point.
Initialized with a small attitude perturbation ($\Delta\theta=1^\circ$,
$\Delta\phi=5^\circ$) --- an offset representative of a minor upset, whose
exact size is immaterial since the divergence is exponential --- the
unstable spiral mode takes over: the bank angle diverges to
${\approx}63^\circ$; the aircraft accelerates and enters a steadily
descending spiral --- the behavior predicted by Table~\ref{tab:poles}.
No other disturbance is applied; the perturbation acts only through the
initial condition.

\begin{figure}[!t]
\centering
\includegraphics[width=\linewidth]{fig_malha_aberta.png}
\caption{Open-loop nonlinear simulation with commands frozen at trim.
Solid: initialized exactly at trim (remains at equilibrium). Dashed: small
initial perturbation --- the unstable spiral mode diverges.}
\label{fig:openloop}
\end{figure}
